\chapter{Vetores e retas no plano}\label{ch:g11:vect}

Os \omterm{def:g11:vect:vector}{vetores} codificam deslocamentos; as retas são suas trajetórias. A
principal ferramenta nova deste capítulo é o \emph{\omterm{def:g11:vect:det}{determinante}}, um único
número que decide se dois \omterm{def:g11:vect:vector}{vetores} são \omterm{def:g11:vect:collinear}{colineares} --- e, a partir daí,
produz \omterm{def:g10:algebra:equation}{equações} de retas e resolve por puro cálculo as questões de
colinearidade e paralelismo. As mesmas ideias se estendem ao espaço no
\cref{ch:g12:space}.

\section{Vetores: uma retomada}

\begin{definition}[Vetor, coordenadas]\label{def:g11:vect:vector}
Um \emph{vetor}\index{vetor} $\vec u$ representa um deslocamento: todas
as setas de mesma direção, mesmo sentido e mesmo comprimento representam
o mesmo vetor. Em um \omterm{def:g10:coordgeom:system}{sistema de coordenadas}, $\vec u\,(x, y)$ registra o
deslocamento ($x$ na horizontal, $y$ na vertical); para pontos
$A(x_A, y_A)$ e $B(x_B, y_B)$,
\[
\vect{AB}\,\bigl(x_B - x_A,\ y_B - y_A\bigr).
\]
Os vetores somam-se coordenada a coordenada, e $\lambda\vec u$ tem
\omterm{def:g10:coordgeom:system}{coordenadas} $(\lambda x, \lambda y)$.
\end{definition}

\begin{proposition}[Ponto médio]\label{prop:g11:vect:midpoint}
O \omterm{prop:g10:coordgeom:midpoint}{ponto médio} $M$ do segmento $[AB]$ tem \omterm{def:g10:coordgeom:system}{coordenadas}
$\left(\frac{x_A + x_B}{2},\ \frac{y_A + y_B}{2}\right)$.
\end{proposition}

\begin{proof}
$M$ é caracterizado por $\vect{AM} = \frac12\vect{AB}$: coordenada a
coordenada, $x_M - x_A = \frac12(x_B - x_A)$, logo
$x_M = \frac{x_A + x_B}{2}$, e analogamente para $y_M$.
\end{proof}

\section{Colinearidade e o determinante}

\begin{definition}[Vetores colineares]\label{def:g11:vect:collinear}
Dois \omterm{def:g11:vect:vector}{vetores} são \emph{colineares}\index{vetores colineares} se um é
múltiplo escalar do outro: $\vec v = \lambda \vec u$ para algum
$\lambda \in \R$ (ou $\vec u = \vec 0$). Geometricamente: eles têm a
mesma direção, ou um deles é nulo.
\end{definition}

\begin{definition}[Determinante]\label{def:g11:vect:det}
O \emph{determinante}\index{determinante} de $\vec u\,(x, y)$ e
$\vec v\,(x', y')$ é
\[
\det(\vec u, \vec v) =
\begin{vmatrix} x & x' \\ y & y' \end{vmatrix}
= x y' - x' y .
\]
\end{definition}

\begin{theorem}[Critério de colinearidade]\label{thm:g11:vect:collinearity}
Dois \omterm{def:g11:vect:vector}{vetores} $\vec u$ e $\vec v$ são \omterm{def:g11:vect:collinear}{colineares} se, e somente se,
$\det(\vec u, \vec v) = 0$.
\end{theorem}

\begin{proof}
($\Rightarrow$) Se $\vec v = \lambda\vec u$, então $x' = \lambda x$ e
$y' = \lambda y$, logo
$\det(\vec u, \vec v) = x(\lambda y) - (\lambda x)y = 0$; se
$\vec u = \vec 0$, o \omterm{def:g11:vect:det}{determinante} é trivialmente $0$.

($\Leftarrow$) Suponha $xy' - x'y = 0$ com $\vec u \neq \vec 0$, digamos
$x \neq 0$ (o caso $y \neq 0$ é simétrico). Ponha
$\lambda = \frac{x'}{x}$. Então $x' = \lambda x$ por construção, e a
relação $xy' = x'y = \lambda x y$ dá, após dividir por $x \neq 0$,
$y' = \lambda y$. Logo $\vec v = \lambda \vec u$.
\end{proof}

\begin{example}\label{ex:g11:vect:collinear}
$\vec u\,(3, -2)$ e $\vec v\,(-6, 4)$:
$\det = 3 \times 4 - (-6) \times (-2) = 12 - 12 = 0$, \omterm{def:g11:vect:collinear}{colineares} (de
fato, $\vec v = -2\vec u$). $\vec u\,(3, -2)$ e $\vec w\,(2, 1)$:
$\det = 3 + 4 = 7 \neq 0$, não são \omterm{def:g11:vect:collinear}{colineares}.
\end{example}

\begin{method}[Colinearidade de três pontos]\label{met:g11:vect:alignment}
$A$, $B$, $C$ são \omterm{def:g11:vect:collinear}{colineares} se, e somente se, $\vect{AB}$ e $\vect{AC}$
são \omterm{def:g11:vect:collinear}{vetores colineares}: calcule os dois \omterm{def:g11:vect:vector}{vetores} e verifique se
$\det\bigl(\vect{AB}, \vect{AC}\bigr) = 0$.
\end{method}

\section{Equações cartesianas de retas}

\begin{definition}[Vetor diretor]\label{def:g11:vect:direction}
Um \emph{vetor diretor}\index{vetor diretor} de uma reta $d$ é qualquer
\omterm{def:g11:vect:vector}{vetor} não nulo $\vec u$ tal que $\vect{AB}$ seja \omterm{def:g11:vect:collinear}{colinear} com $\vec u$
para todos os pontos $A, B$ de $d$.
\end{definition}

\begin{theorem}[Equação cartesiana]\label{thm:g11:vect:cartesian}
Uma reta de \omterm{def:g11:vect:direction}{vetor diretor} $\vec u\,(-b, a)$ admite uma \omterm{def:g10:algebra:equation}{equação} da forma
\[
ax + by + c = 0
\qquad (a, b) \neq (0, 0).
\]
Reciprocamente, toda \omterm{def:g10:algebra:equation}{equação} desse tipo descreve uma reta de \omterm{def:g11:vect:direction}{vetor
diretor} $\vec u\,(-b, a)$.
\end{theorem}

\begin{proof}
Seja $d$ passando por $A(x_A, y_A)$ com direção $\vec u\,(-b, a)$. Um
ponto $M(x, y)$ está em $d$ exatamente quando $\vect{AM}$ é \omterm{def:g11:vect:collinear}{colinear} com
$\vec u$, isto é (\cref{thm:g11:vect:collinearity}),
\[
0 = \det\bigl(\vect{AM}, \vec u\bigr)
= (x - x_A)\,a - (-b)(y - y_A)
= ax + by - (a x_A + b y_A),
\]
uma \omterm{def:g10:algebra:equation}{equação} da forma anunciada com $c = -(a x_A + b y_A)$.
Reciprocamente, as soluções de $ax + by + c = 0$ com, digamos,
$b \neq 0$ são os pontos $\left(x, -\frac{a x + c}{b}\right)$: subtrair
duas delas mostra que toda corda é \omterm{def:g11:vect:collinear}{colinear} com $(-b, a)$, e sempre
existe uma solução --- trata-se da reta que passa por ela e é dirigida
por $(-b, a)$.
\end{proof}

\begin{method}[Reta por dois pontos]\label{met:g11:vect:twopoints}
Para achar uma \omterm{def:g10:algebra:equation}{equação} da reta $(AB)$: calcule o \omterm{def:g11:vect:direction}{vetor diretor}
$\vect{AB}\,(x_B - x_A, y_B - y_A)$; identifique-o com $(-b, a)$, isto é,
tome $a = y_B - y_A$ e $b = -(x_B - x_A)$; determine $c$ substituindo as
\omterm{def:g10:coordgeom:system}{coordenadas} de $A$.
\end{method}

\begin{example}\label{ex:g11:vect:twopoints}
Reta por $A(1, 2)$ e $B(4, 3)$: $\vect{AB}\,(3, 1)$, logo $a = 1$,
$b = -3$, e a \omterm{def:g10:algebra:equation}{equação} é $x - 3y + c = 0$. Substituindo $A$:
$1 - 6 + c = 0$, logo $c = 5$. \omterm{def:g10:algebra:equation}{Equação}: $x - 3y + 5 = 0$. Confira com
$B$: $4 - 9 + 5 = 0$.
\end{example}

\begin{remark}
Quando $b \neq 0$, a \omterm{def:g10:algebra:equation}{equação} pode ser resolvida em $y$: $y = mx + p$, com
\omterm{def:g10:reffunc:affine}{coeficiente angular} $m = -\frac ab$. A forma cartesiana
$ax + by + c = 0$ é mais geral: ela cobre também as retas verticais
($b = 0$), que não têm \omterm{def:g10:reffunc:affine}{coeficiente angular}.
\end{remark}

\begin{omfigure}
\begin{tikzpicture}
\begin{axis}[omaxis,
  width=8.6cm, height=5.8cm,
  xmin=-1.6, xmax=5.6, ymin=-0.9, ymax=4.4,
  xtick={1,4}, ytick={2,3}]
  \addplot[omDef, thick, domain=-1.4:5.4] {(x+5)/3};
  \fill (axis cs:1,2) circle (1.5pt)
    node[above left, font=\small] {$A$};
  \fill (axis cs:4,3) circle (1.5pt)
    node[above left, font=\small] {$B$};
  \draw[-Stealth, omThm, thick] (axis cs:1,2) -- (axis cs:2.5,2.5)
    node[midway, below right, font=\small] {$\vec u$};
\end{axis}
\end{tikzpicture}

{\small A reta $x - 3y + 5 = 0$ do \cref{ex:g11:vect:twopoints}, seu
\omterm{def:g11:vect:direction}{vetor diretor} $\vec u = \vect{AB}$ (em vermelho, aqui reduzido pela
metade) e os dois pontos usados para montar a \omterm{def:g10:algebra:equation}{equação}.}
\end{omfigure}

\section{Representação paramétrica}

\begin{definition}[Representação paramétrica]\label{def:g11:vect:parametric}
A reta que passa por $A(x_A, y_A)$ com direção $\vec u\,(\alpha, \beta)$
é o conjunto dos pontos
\[
M\bigl(x_A + t\alpha,\ y_A + t\beta\bigr), \qquad t \in \R .
\]
O número $t$ é o \emph{parâmetro}: cada valor de $t$ produz um ponto da
reta, como se a percorrêssemos com velocidade constante $\vec u$.
\end{definition}

\begin{example}
A reta do \cref{ex:g11:vect:twopoints} também é
$\{(1 + 3t,\ 2 + t) : t \in \R\}$: $t = 0$ dá $A$ e $t = 1$ dá $B$.
Eliminar $t$ ($t = y - 2$, logo $x = 1 + 3(y-2)$) recupera
$x - 3y + 5 = 0$.
\end{example}

\section{Posições relativas de duas retas}

\begin{proposition}[Paralelas ou concorrentes]\label{prop:g11:vect:positions}
Duas retas de \omterm{def:g11:vect:direction}{vetores diretores} $\vec u$ e $\vec v$ são paralelas se, e
somente se, $\det(\vec u, \vec v) = 0$. Em forma de \omterm{def:g10:algebra:equation}{equação}:
$d : ax + by + c = 0$ e $d' : a'x + b'y + c' = 0$ são paralelas se, e
somente se, $ab' - a'b = 0$. Retas não paralelas se encontram em
exatamente um ponto.
\end{proposition}

\begin{proof}
Paralelismo significa direções \omterm{def:g11:vect:collinear}{colineares}, o que é a
\cref{thm:g11:vect:collinearity}; com $\vec u\,(-b, a)$ e
$\vec v\,(-b', a')$, o \omterm{def:g11:vect:det}{determinante} é $(-b)a' - (-b')a = ab' - a'b$. Se
as retas não são paralelas, resolver as duas \omterm{def:g10:algebra:equation}{equações} simultaneamente
(por substituição ou combinação) leva a uma única solução: o sistema se
comporta como um sistema $2 \times 2$ de \omterm{def:g11:vect:det}{determinante} não nulo.
\end{proof}

\begin{method}[Interseção de duas retas]\label{met:g11:vect:intersection}
Resolva o sistema das duas \omterm{def:g10:algebra:equation}{equações}: isole uma incógnita em uma delas e
substitua na outra, ou combine as \omterm{def:g10:algebra:equation}{equações} para eliminar uma incógnita.
Confira sempre o ponto obtido nas \emph{duas} \omterm{def:g10:algebra:equation}{equações}.
\end{method}

\begin{example}\label{ex:g11:vect:intersection}
Intersecte $d: x - 3y + 5 = 0$ e $d': 2x + y - 4 = 0$. De $d'$:
$y = 4 - 2x$. Substituindo em $d$:
\[
x - 3(4 - 2x) + 5 = 0
\iff x - 12 + 6x + 5 = 0
\iff 7x = 7
\iff x = 1,
\]
e então $y = 2$. As retas se encontram em $(1, 2)$. Verificação em $d$:
$1 - 6 + 5 = 0$.
\end{example}

\section{Exercícios}

\begin{exercise}[$\star$]\label{exo:g11:vect:1}
Dados $A(2, -1)$, $B(5, 3)$ e $C(-1, 1)$, calcule as \omterm{def:g10:coordgeom:system}{coordenadas} de
$\vect{AB}$, $\vect{AC}$, $\vect{AB} + \vect{AC}$ e
$2\vect{AB} - 3\vect{AC}$, e o \omterm{prop:g10:coordgeom:midpoint}{ponto médio} de $[BC]$.
\end{exercise}

\begin{exercise}[$\star$]\label{exo:g11:vect:2}
Os \omterm{def:g11:vect:vector}{vetores} são \omterm{def:g11:vect:collinear}{colineares}?
\[
\vec u\,(4, -6) \text{ e } \vec v\,(-2, 3); \qquad
\vec u\,(2, 5) \text{ e } \vec v\,(3, 7); \qquad
\vec u\,(0, 3) \text{ e } \vec v\,(0, -8).
\]
\end{exercise}

\begin{exercise}[$\star$]\label{exo:g11:vect:3}
Encontre uma \omterm{def:g10:algebra:equation}{equação} cartesiana da reta que passa por $A(2, 1)$ com \omterm{def:g11:vect:direction}{vetor
diretor} $\vec u\,(3, -1)$ e, em seguida, da reta que passa por $B(0, 4)$
e $C(2, 0)$.
\end{exercise}

\begin{exercise}[$\star$]\label{exo:g11:vect:4}
Para a reta $d: 3x - 2y + 6 = 0$, dê um \omterm{def:g11:vect:direction}{vetor diretor}, as \omterm{def:g10:numbers:interunion}{interseções} com
os dois eixos coordenados e decida se os pontos $P(2, 6)$ e $Q(1, 4)$
pertencem a $d$.
\end{exercise}

\begin{exercise}[$\star\star$]\label{exo:g11:vect:5}
Os pontos $A(1, 2)$, $B(4, 8)$, $C(-2, -4)$ são \omterm{def:g11:vect:collinear}{colineares}? Mesma
pergunta para $D(0, 1)$, $E(2, 4)$, $F(5, 9)$.
\end{exercise}

\begin{exercise}[$\star\star$]\label{exo:g11:vect:6}
Determine o ponto de \omterm{def:g10:numbers:interunion}{interseção}, se houver:
\[
d: 2x + y - 7 = 0 \ \text{ e }\ d': x - y + 1 = 0;
\qquad
e: x - 2y + 3 = 0 \ \text{ e }\ e': -2x + 4y + 1 = 0 .
\]
\end{exercise}

\begin{exercise}[$\star\star$]\label{exo:g11:vect:7}
Dê uma representação paramétrica da reta $d: 2x - y + 3 = 0$ e uma
\omterm{def:g10:algebra:equation}{equação} cartesiana da reta
$\bigl\{(1 + 2t,\ -3 + t) : t \in \R\bigr\}$.
\end{exercise}

\begin{exercise}[$\star\star$]\label{exo:g11:vect:8}
Encontre a \omterm{def:g10:algebra:equation}{equação} da reta que passa por $P(1, -2)$ e é paralela a
$d: 4x - 3y + 2 = 0$, e o valor de $m$ para o qual a reta
$mx + 2y - 5 = 0$ é paralela a $d$.
\end{exercise}

\begin{exercise}[$\star\star$]\label{exo:g11:vect:9}
Sejam $A(-1, 0)$, $B(3, 2)$ e $C(1, -4)$. Encontre uma \omterm{def:g10:algebra:equation}{equação} cartesiana
da \emph{\omterm{def:g10:stats:median}{mediana}} do triângulo $ABC$ que parte de $C$ (a reta que liga
$C$ ao \omterm{prop:g10:coordgeom:midpoint}{ponto médio} de $[AB]$).
\end{exercise}

\begin{exercise}[$\star\star$]\label{exo:g11:vect:10}
Para que valor (ou valores) do parâmetro $m$ os \omterm{def:g11:vect:vector}{vetores} $\vec u\,(m, 2)$
e $\vec v\,(3, m - 1)$ são \omterm{def:g11:vect:collinear}{colineares}?
\end{exercise}

\begin{exercise}[$\star\star\star$]\label{exo:g11:vect:11}
Seja $ABCD$ um paralelogramo, $I$ o \omterm{prop:g10:coordgeom:midpoint}{ponto médio} de $[AB]$ e $J$ o ponto
definido por $\vect{DJ} = \frac23 \vect{DI}$. Trabalhando no \omterm{def:g10:coordgeom:system}{sistema de
coordenadas} em que $A(0,0)$, $B(1,0)$, $D(0,1)$ (de modo que $C(1,1)$):
\begin{enumerate}
  \item Calcule as \omterm{def:g10:coordgeom:system}{coordenadas} de $I$ e $J$.
  \item Mostre que $A$, $J$ e $C$ são \omterm{def:g11:vect:collinear}{colineares}. (Um fato bonito: a reta
        $(DI)$ corta a diagonal $(AC)$ a um terço do caminho.)
\end{enumerate}
\end{exercise}

\section{Problema: Rotas de colisão e coordenadas espertas}

\begin{problem}[{Problema de fim de semana --- cruzar trajetórias não
é colidir: a aproximação mínima no mar e teoremas demonstrados
escolhendo o referencial certo}]
\label{pb:g11:vect:1}
As rotas retas de dois navios se cruzam --- os navios têm de
colidir? Só se chegarem ao cruzamento \emph{ao mesmo tempo}:
trajetórias são conjuntos de pontos, movimentos são paramétricos,
e confundir as duas coisas já afundou navios de verdade. Este
problema opera um pequeno serviço de controle de tráfego marítimo
com retas paramétricas (\cref{def:g11:vect:parametric}) e depois
volta-se para a geometria pura com o outro superpoder das
\omterm{def:g10:coordgeom:system}{coordenadas}: \emph{escolhê-las} com esperteza, como no
\cref{exo:g11:vect:11}, para que os teoremas clássicos se
demonstrem sozinhos.

\medskip
\noindent\textbf{Parte I --- Retas, de três maneiras.}
\begin{enumerate}
  \item Dê uma representação paramétrica da reta que passa por
        $A(1, 2)$ com \omterm{def:g11:vect:direction}{vetor diretor} $\vec u(3, -1)$ e depois
        elimine o parâmetro para obter uma \omterm{def:g10:algebra:equation}{equação} cartesiana.
  \item Para a reta $2x - 5y + 3 = 0$: leia um \omterm{def:g11:vect:direction}{vetor diretor}
        (\cref{thm:g11:vect:cartesian}), encontre um ponto e
        escreva uma representação paramétrica.
  \item Calcule o ponto de \omterm{def:g10:numbers:interunion}{interseção} das duas retas das questões
        1 e 2 (\cref{met:g11:vect:intersection}).
  \item Confirme com o \omterm{def:g11:vect:det}{determinante}
        (\cref{thm:g11:vect:collinearity}) que essas duas retas
        tinham mesmo de se cruzar; mostre em seguida que
        $2x - 5y + 3 = 0$ e $-4x + 10y + 1 = 0$ são estritamente
        paralelas.
  \item $A(2, 1)$, $B(5, 3)$, $C(11, 7)$ são \omterm{def:g11:vect:collinear}{colineares}
        (\cref{met:g11:vect:alignment})?
\end{enumerate}

\medskip
\noindent\textbf{Parte II --- O controle de tráfego marítimo.}
No instante $t$ (em horas), o navio $P$ está em $(2t,\ 3 + t)$ e o
navio $Q$ em $(10 - t,\ 2t - 1)$ (distâncias em milhas náuticas).
\begin{enumerate}[resume]
  \item Encontre as \omterm{def:g10:algebra:equation}{equações} cartesianas das duas
        \emph{trajetórias} e calcule onde elas se cruzam.
  \item Em que instante cada navio passa pelo ponto de
        cruzamento? Os navios colidem? Enuncie o aviso geral em
        uma frase: trajetórias que se cruzam contra movimentos
        que se encontram.
  \item Calcule o quadrado da distância $D^2(t)$ entre os navios
        no instante $t$ e reduza-o a um \omterm{def:g11:quad:quadratic}{trinômio do segundo grau}
        em $t$. Em que instante os navios ficam mais próximos, e
        a que distância chegam a estar (a tecnologia do vértice
        do \cref{pb:g11:quad:1})?
  \item O regulamento exige uma separação de pelo menos $1$
        milha. Ela é violada? Se for, resolva $D^2(t) < 1$ e dê a
        duração da violação.
  \item Explique por que, para \emph{quaisquer} dois navios em
        rotas retas com velocidades constantes, $D^2(t)$ é sempre
        uma \omterm{def:g11:func:function}{função} do segundo grau em $t$ --- de modo que o
        ``ponto de \omterm{def:g10:numbers:approx}{aproximação} mínima'' é sempre um cálculo de
        vértice.
  \item Interceptação: uma lancha de patrulha parte da origem em
        $t = 0$ com velocidade constante $(a, b)$ e deve
        encontrar o navio $P$ exatamente em $t = 3$. Calcule o
        $(a, b)$ necessário e a velocidade da lancha.
\end{enumerate}

\medskip
\noindent\textbf{Parte III --- \omterm{def:g10:coordgeom:system}{Coordenadas} espertas.}
Trabalhe no referencial do \cref{exo:g11:vect:11}: o
paralelogramo $ABCD$ vira $A(0,0)$, $B(1,0)$, $C(1,1)$,
$D(0,1)$; seja $I$ o \omterm{prop:g10:coordgeom:midpoint}{ponto médio} de $[AB]$ e $K$ o \omterm{prop:g10:coordgeom:midpoint}{ponto médio} de
$[CD]$.
\begin{enumerate}[resume]
  \item Calcule uma \omterm{def:g10:algebra:equation}{equação} cartesiana da reta $(DI)$.
  \item Intersecte $(DI)$ com a diagonal $(AC)$: recupere o fato
        do \cref{exo:g11:vect:11} --- o corte fica exatamente a
        um terço do caminho ao longo da diagonal.
  \item Agora a reta $(BK)$: encontre sua \omterm{def:g10:algebra:equation}{equação}, intersecte-a
        com $(AC)$ e conclua o teorema clássico completo: \emph{as
        cevianas $(DI)$ e $(BK)$ dividem a diagonal $(AC)$ em
        três partes iguais}.
  \item Por que foi legítimo demonstrar o teorema nesse
        referencial particular? (O que a escolha de \omterm{def:g10:coordgeom:system}{coordenadas}
        preserva e do que o enunciado trata?) Responda em duas ou
        três frases.
  \item Mais uma trissecção para praticar: seja $E$ o \omterm{prop:g10:coordgeom:midpoint}{ponto médio}
        de $[BC]$. Onde a reta $(AE)$ corta a \emph{outra}
        diagonal, $(DB)$?
\end{enumerate}

\medskip
\noindent\textbf{Parte IV --- Enigmas de posição.}
\begin{enumerate}[resume]
  \item As três retas $x + y = 4$, \ $2x - y = 2$, \
        $x - 2y = -2$ são concorrentes? (Intersecte duas e teste
        a terceira.)
  \item Para cada real $m$, seja $d_m$ a reta
        $mx - y + 2 - 3m = 0$. Mostre que \emph{toda} $d_m$ passa
        por um mesmo \omterm{pb:g10:functions:1}{ponto fixo}. (Agrupe os termos em $m$.)
  \item Na família $(d_m)$: qual membro é paralelo a
        $y = 2x + 7$? Qual passa pela origem?
  \item Final --- os três trajes de uma reta: cartesiano (teste
        de pertinência em uma substituição), reduzido
        (\omterm{def:g10:reffunc:affine}{coeficiente angular} e \omterm{def:g10:functions:graph}{gráfico} à vista), paramétrico
        (movimento, horários, interceptação); o \omterm{def:g11:vect:det}{determinante} como
        oráculo do paralelismo; e o referencial bem escolhido
        como fábrica de teoremas. Uma frase para cada.
\end{enumerate}
\end{problem}
